\paragraph{Conjugate-coordinate class times a fuzzy factor: the Hessian of a tree stabiliser on \(K\), the one-loop weights on the cubic branch and the exchange of the scale of \(K\) (conditional; bounded no-go; correction of the cubic and mass branches of the preceding paragraph).}
This paragraph supplements the preceding one, in which a tree cubic term and a tree mass term on \(K\) were added to the tree term \(V_{\mathrm t}\) while \(U_4\), \(U_2\) and \(Y\) were those of the operators \(H_0+\Sigma^{AB}\mathbb F_{AB}\) of the paragraph on product backgrounds. Keep its setting and notation, with \(K=S^2_N\), \(X^i=\chi J^i\) in three of the six directions, and take the cubic term of Equation~\eqref{eq:Smatrix} on those three directions with \(C_{ijk}=\mathrm i\,c_3\,\epsilon_{ijk}/g^2\), \(c_3\) real, so that the \(K\) part of the tree action is \(g^{-2}\bigl[-\tfrac14\operatorname{tr}[X^i,X^j]^2+\tfrac{\mathrm ic_3}{3}\epsilon_{ijk}\operatorname{tr}X^iX^jX^k\bigr]\), which is real for Hermitian matrices (for real \(C_{ijk}\) the cubic trace is imaginary). Definition~\ref{def:matrix-coords-ext} does not extend the cubic term to the transverse coordinates, so this term is a supplied input, as in the preceding paragraph. Write \(r_3=c_3/2\chi\) and \(\kappa_3=2-c_3/\chi=2(1-r_3)\); \(\kappa_3\) is not the weight \(\kappa\) of the paragraph on fibre-point forms.

\emph{Hessian.} On \(S^2_N\) the tree term is \(V_{\mathrm t}=\bar a(\chi^4-\tfrac23c_3\chi^3)\), stationary at \(\chi=c_3/2\) with \(\partial^2V_{\mathrm t}/\partial(\ln\chi)^2=4\bar a\chi^4>0\) there; this is the fuzzy-sphere solution of the matrix model with a cubic term~\cite{IsoKimuraTanakaWakatsuki2001}. The second variation of the cubic term is \(\mathrm ic_3\,\epsilon_{ijk}\operatorname{tr}(a^ia^jX^k)\): it acts on the three components of the vector along \(K\) only and is linear in the adjoint action of \(X^k\). With \(L^2\) the adjoint Casimir and \(S\) the spin-one generators on the index \(i\), the operator of those three components in the background Feynman gauge is exactly
\[
  H_V^K=\chi^2\bigl(L^2-\kappa_3\,S\cdot L\bigr),\qquad \kappa_3=2-\frac{c_3}\chi ,
\]
with eigenvalues \(\chi^2(l(l+1)-\kappa_3l)\), \(\chi^2(l(l+1)+\kappa_3)\), \(\chi^2(l(l+1)+\kappa_3(l+1))\) and multiplicities \(2l+3\), \(2l+1\), \(2l-1\) (Hessian of the action built from the matrices at \(N\le5\), six pairs \((\chi,c_3)\), residual below \(10^{-13}\)); \(\kappa_3=2\) is the spectrum of the preceding paragraph. The ghost operator, the spinor operator and the other seven components of the vector are unchanged, since the cubic term contains neither ghosts nor spinors. There is no negative mode for \(-1\le\kappa_3\le2\), that is \(0\le c_3\le3\chi\); at \(\kappa_3=-2\) the mode \(l=1\), \(j=0\) has the eigenvalue \(-2\chi^2\). At \(\chi=c_3/2\) the operator is \(\chi^2L^2\) on the three components, the index part \(v^K\) vanishes identically, and the Hessian without gauge fixing annihilates the gauge directions, which it does not at \(c_3=0\). A tree mass term \(\tfrac12\mu^2\operatorname{tr}X^iX_i\) on \(K\) adds \(\mu^2\) to the same three components and to nothing else.

\emph{Weights with the Hessian.} The factorisation of the paragraph on product backgrounds is an identity for commuting operators of the two factors and holds with \(v^K\) the index part of \(H_V^K\). Its statement that the index parts of the vector and of the spinors agree at second order in the flux does not: with weights \(w\) and eigenvalues \(\lambda\) of \(v^K-\tau^K\), \(\sum w=\sum w\lambda=0\) and, exactly,
\[
  \sum w\lambda^2=(\kappa_3^2-4)\,\chi^4N^2(N^2-1)=(c_3^2\chi^2-4c_3\chi^3)\,N^2(N^2-1),
\]
so \(s_0^M(v^K-\tau^K)\) starts at second order in the flux of \(K\). Then \(\hat U\) of the preceding paragraph does not converge at small proper time. With a lower limit \(s_\Lambda\) the term \(U_4\) is replaced by \(-(F/32\pi^2)\sqrt{|g|}\,\mathcal U\), where, with \(E_3\) the exponential integral of order three,
\[
  \mathcal U=\int_{s_\Lambda}^\infty\frac{\mathrm ds}{s^3}\bigl(v^K-\tau^K\bigr)
  =\frac1{s_\Lambda^2}\sum wE_3(s_\Lambda\lambda)
  =-\tfrac12\sum w\lambda^2\ln\lambda+\tfrac12\Bigl(\sum w\lambda^2\Bigr)\Bigl(\ln\frac1{s_\Lambda}-\gamma_{\mathrm E}+\tfrac32\Bigr)+O(s_\Lambda\lambda_{\max}),
\]
whose logarithm multiplies a term in \(\chi^3\), of the form of the cubic term, and a term in \(\chi^2\), of the form of a mass term on \(K\), which Equation~\eqref{eq:Smatrix} does not contain; both carry \(\sqrt{|g|}\), not \(\rho_M\). So in the Poisson limit on \(M\) the cubic branch is not closed at one loop without a mass counterterm, and the scaling argument of the preceding paragraph (dependence on \(\chi^2/\lambda\) alone) does not apply at fixed \(c_3\). The remainder of the expansion is of order one where \(s_\Lambda\chi^2N^2\) is, so the values below are those of the exact form. At fixed \(s_\Lambda\chi^2N^2\) the regulated \(\mathcal U\) is of order \(N^4\), against \(N^2\) at \(c_3=0\), and it is negative at every sampled limit \(10^{-4}\le s_\Lambda\chi^2N^2\le30\) for every sampled \(-1<\kappa_3\le1.9\) (\(N=16\), \(64\), \(512\)); for \(\kappa_3\) closer to \(2\) it is positive at some sampled limits at \(N=16\). Where it is negative, \(V_0>0\) and \(q_b^2<0\) in the vacuum term of the preceding paragraph, the sign for which that paragraph finds a Yukawa pole, but the value depends on the regulator. The weights with two derivatives stay finite: \(T\) is unchanged, and \(\hat V/\hat T\to\kappa_3^2/4\) (exact spectra, \(N\le512\)), so that at large \(N\)
\[
  \frac VT=(1-r_3)^2,\qquad W_R=V-T=-T\,r_3(2-r_3),
\]
with \(V=0\) and \(W_R=-T\) exactly at \(r_3=1\). There, in the units of the paragraph on product backgrounds and at equal weights, \(c_T=\tfrac16W_R-4W_X=\tfrac13T>0\), and \(\tfrac16W_R-2W_X=\tfrac1{12}T\) in the presentation \(r=-2\) of that paragraph; the forms of the paragraph on fibre-point forms, which take \(V>0\), degenerate, since \(A+B=\tfrac18V\to0\), \(\varrho=(V+2T)/3V\to\infty\) and \(c^{\mathrm{TT}}\to0\). At large \(N\) the condition \(V=T\) holds only at \(r_3=0\) and at \(r_3=2\), and at \(r_3=2\) the vector operator has the negative mode above. For the mass term \(\sum w\lambda=3\mu^2N^2\ne0\), so \(W_R\) itself is logarithmic at small proper time.

\emph{The example of the preceding paragraph.} The forms of \(\beta_K\) of the preceding paragraph with a tree cubic or a tree mass term, the value \(\beta_K=1\) in the cubic case and at its example at \(x=\tfrac3{10}\), and its availability at negligible \(Y\) for \(\hat U<0\), use \(U_4\), \(U_2\) and \(Y\) of the operators at \(c_3=\mu=0\); they are the one-loop forms of the action with that term only as \(c_3/\chi\to0\) and \(\mu\to0\). The example has \(r_3=1.168\), \(\kappa_3=-0.336\), so with the Hessian \(V/T=0.028\) and \(\varrho=23.9\) there. In the units of that example the term that replaces \(U_4=-0.09\) is \(41.1\) at \(N=16\) and \(659\) at \(N=64\) for \(s_\Lambda\chi^2N^2=\tfrac12\), against the tree quartic \(7.935\). At large \(N\), \(U_2\) becomes \(\bigl[2(4-\kappa_3^2)\vartheta^{-1}-\tfrac34(8-\kappa_3^2)\bigr]x\), which at \(\vartheta=1\), its value at large \(\tau\), is \(0.56\) in place of \(-0.9\), with derivative \(1.70\) with respect to \(\ln\chi\) at fixed \(c_3\) in place of \(-1.8\), and \(Y\) is unchanged, since it comes from the scalar trace on \(K\); so at \(\vartheta=1\) the one-loop slope without the term in \(\mathcal U\) is \(9.2\), larger than \(5.34\), the size of the tree slope there (and it stays larger for the sampled \(1\le\vartheta\le2\)), and the slope of the term in \(\mathcal U\) is positive at every sampled limit \(10^{-4}\le s_\Lambda\chi^2N^2\le30\) (\(N=16\), \(64\), \(512\)). The example is therefore not a stationary point of the one-loop potential with that cubic term for these lower limits without counterterms; with counterterms of the cubic and mass forms, or with a regulator derived from the matrices of \(M\), this is not determined. The relation \(\beta_K=-2\,\partial_{\ln\chi}(\rho_MV_{\mathrm t}/g^2)/U_{\chi\chi}\) holds where every one-loop term responds to the wave through \(\sqrt{|g|}\) alone, which the preceding paragraph shows for \(U_4\), the index part of \(U_2\) and the fibre-constant part of \(Y\); for \(V\ne T\) the term of \(U_2\) proportional to \(W_R\) enters at leading order and its response is not computed, so on the cubic branch the relation is conditional on it.

\emph{Tree curve and the ratio of the one-loop to the tree terms.} With \(\tilde U=\rho_M\bar a\chi^4/g^2\), \(U'_{\mathrm{loop}}\) the derivative of the one-loop density with respect to \(\ln\chi\) and \(\delta\) its second derivative divided by \(4\tilde U\), stationarity gives, for the quartic and cubic terms,
\[
  1-r_3=-\frac{U'_{\mathrm{loop}}}{4\tilde U},\qquad \frac VT=\Bigl(\frac{U'_{\mathrm{loop}}}{4\tilde U}\Bigr)^2 ,
\]
the second at large \(N\), and, under the condition just stated, \(\beta_K=2(r_3-1)/(4-3r_3+\delta)\). So on this branch \(V=T\) requires a one-loop slope equal in size to the slope of the tree quartic, and, under the same condition, \(\beta_K=1\) requires \(r_3=(6+\delta)/5\), where \(V/T=(1+\delta)^2/25\), and at \(r_3=2\), \(\beta_K=2/(\delta-2)\). Where the one-loop term is small against the tree term the minimum is at \(r_3\to1\), which, unlike the quartic term alone, exists at tree level; there \(V/T\to0\) and, under that condition, \(\beta_K\to0\), the case (S2). Within one loop, on \(S^2_N\) with this cubic term at \(c_3\ne0\) and at large \(N\), no point has \(V=T\) and no negative mode, so none has \(\beta_K=1\) with \(V=T\). At the example of the preceding paragraph, with its operators, the ratios of the one-loop to the tree contributions are \(-1\) for the slope (by stationarity; the one-loop slope is \(0.168\) of the slope of the quartic) and \(-5.04/15.72=-0.32\) for \(U_{\chi\chi}\), and the vacuum term has no tree part among the computed forms. These are values at one loop; where the ratios are of order one they are not controlled, and higher loops are not computed.

\emph{Exchange of the scale of \(K\).} For a real wave with \(\delta\ln\sqrt{|g|}=\sigma_1=-t/2m^2\) at first order, as in the static channel of the preceding paragraph, the terms without derivatives are \(\tfrac12U_{\chi\chi}(\delta\ln\chi)^2+U'_{\mathrm{loop}}\,\delta\ln\chi\,\sigma_1+U_4\sigma_1^2\) plus the second-order parts of \(U_2\) and of the first term, with the first-order responses of the preceding paragraph. Eliminating \(\delta\ln\chi\) for \(q^2\) below \(m_\chi^2\) gives \((U_4-\tfrac14\beta_KU'_{\mathrm{loop}})\sigma_1^2\), with \(\beta_K=2U'_{\mathrm{loop}}/U_{\chi\chi}\), so the term \(q_*^2\) of the preceding paragraph contains
\[
  q_*^2\big|_\chi=-\frac{\beta_KU'_{\mathrm{loop}}L}{8m^4}=\frac{\beta_KU'_{\mathrm{loop}}}{4U_4}\,q_b^2 ,
\]
which is non-positive at a local minimum. At the example of the preceding paragraph, with its operators, it is \(-14.8\,q_b^2\), so there the exchange of \(\chi\) exceeds the vacuum term and has the tachyonic sign. The second-order parts of \(U_2\) and of the first term, and the constant of the first-term logarithm, are not computed, so the sign of \(q_*^2-q_b^2\) is not determined.

These statements are for the hypotheses of the preceding paragraph, \(K=S^2_N\) with fixed \(\Lambda^i=J^i\), the cubic term with \(C_{ijk}\propto\mathrm i\epsilon_{ijk}\) on the three directions of \(K\) and no fermion bilinear, the background Feynman gauge away from \(\chi=c_3/2\), one loop, the Poisson limit on \(M\) with a lower proper-time limit and no counterterms where one is stated, and large \(N\) for the ratios of weights and for \(U_2\). The Hessian, the spectra, the moments, the form of \(\mathcal U\) and the relations on the tree curve are exact for these forms; the limits of the weights, the sign of the regulated \(\mathcal U\) and of its slope are numerical at the sampled limits. They do not compute the one-loop potential of \(\chi\) with a regulator derived from the matrices of \(M\), a completion of the cubic term with fermion bilinears or mass terms, other \(K\), the rotating class, the response of the term of \(U_2\) proportional to \(W_R\), the second-order response of \(U_2\) and of the first term, or higher loops; they do not construct the \(K\) selected by the theory, determine \(\mathcal M\), identify \(\bar h\), \(h_\beta\) or \(t\) with quantities of the emergent metric of Chapter~\ref{ch:vierbein-emergence}, discharge Assumption~\ref{ass:gap-psi}, or bear on Assumption~\ref{ass:lorentzian}, so no value or sign of \(G_{\mathrm{eff}}\) and no statement on a long-range force follow. Problems~C, F and~N remain open and PROGRAMME-CONDITIONAL.

